\section{Methodology}
\label{sec:method}

% In this section, we detail our \ourmethod framework, a semantic guidance method that leverages a vision-language model (VLM) as a semantic planner to enhance the instruction grounding capacity of world-action models.
% We first give an overview of the overall formulation.
% We then present its two core components, the semantic planner and the
% semantic-guided world-action model, and finally the three-stage training
% paradigm.
\subsection{Overview}
\label{sec:method_overview}
As shown in Figure~\ref{fig:overview}, the VLM planner takes the current observation $\mathbf{o}_t$ and the language
instruction $\ell$, together with a set of learnable query tokens.
From the hidden states of these query tokens, it predicts two dense feature maps for future keyframes: $\hat{\mathbf{F}}_{\mathrm{tg}}$ grounds the instruction on the task-relevant objects, and $\hat{\mathbf{F}}_{\mathrm{spa}}$ captures the
scene geometry required for precise manipulation.
During training, the two maps are supervised by frozen teacher encoders that
provide target features from future frames, driving the query states to
anticipate how the scene should evolve.

The world-action model, which couples a video expert with an action expert, generates the future video $\mathbf{o}_{t+1:t+T}$ of $T$ frames and the action chunk $\mathbf{a}_{t:t+H-1}$ of $H$ steps under the guidance of this foresight. We model the resulting conditional distribution as
\begin{equation}
    p_{\theta}\!\left(
        \mathbf{a}_{t:t+H-1},\,
        \mathbf{o}_{t+1:t+T}
        \,\middle|\,
        \mathbf{o}_t,\ell,\mathbf{s}_t,
        \hat{\mathbf{F}}_{\mathrm{tg}},
        \hat{\mathbf{F}}_{\mathrm{spa}}
    \right),
    \label{eq:wam_overview}
\end{equation}
where $\theta$ denotes the model parameters and $\mathbf{s}_t$ is the robot state that the action expert receives.

We train the model in three stages: we first train the semantic planner alone to
predict the semantic guidance of future keyframes, then co-train it with the
video expert, and finally add the action expert.
At inference, we discard the teacher encoders, and a single planner forward pass supplies
the foresight that guides generation.

\subsection{Semantic Foresight Learning}
\label{sec:method_planner}

We instantiate the semantic planner with Qwen3.5~\cite{qwen35}.
To predict future features in a single forward pass, we append learnable query
tokens to the image--instruction sequence and take their hidden states.
To capture the semantic knowledge shared by the text-grounded and spatial-aware
targets, we split the tokens into a base group and two modality-specific groups.
Both branches read the base group, while each modality-specific group captures
only what its own target requires.
Branch $m$ then reads the concatenated states
$\mathbf{H}_m=[\mathbf{H}_{\mathrm{base}};\mathbf{H}_m^{\mathrm{spec}}]$.
We reuse the same tokens for $K$ future keyframes at increasing offsets, so the
foresight covers a horizon instead of a single instant.
The planner is applied to each camera view independently and with shared
weights, so the same instruction is grounded separately in every view.
A lightweight resampler $\mathcal{R}_m$ then decodes each $\mathbf{H}_m$ into a
dense feature map.
To produce this map, a bank of $P$ learnable grid queries $\mathbf{U}_m$
cross-attends to $\mathbf{H}_m$ and the image states $\mathbf{E}$.
By attending to both sources, the resampler recovers fine spatial details from the observation while the query states specify how it should change:
\begin{equation}
    \revcolor
    \hat{\mathbf{F}}_m^{(k,v)}
    =
    \mathcal{R}_m\!\left([\mathbf{E}^{(v)};\mathbf{H}_m^{(k,v)}],\,\mathbf{U}_m\right)
    \in\mathbb{R}^{P\times d_m},
    \label{eq:resampler}
\end{equation}
where $m\in\{\mathrm{tg},\mathrm{spa}\}$ indexes the two semantic branches\rev{, $k\in\{1,\dots,K\}$ indexes the future keyframe, and $v\in\{1,\dots,V\}$ indexes the camera view; $\mathbf{E}^{(v)}$ are the image states of view $v$, and $\mathbf{H}_m^{(k,v)}$ are the query states read out at the positions of keyframe $k$ in the sequence of view $v$. The query embeddings and $\mathbf{U}_m$ are shared across $k$ and $v$}.
Both maps use $P{=}256$ tokens to match the token layout of the two teachers, and the feature dimension $d_m$ is $1024$ for the
text-grounded map and $2048$ for the spatial-aware map.

Two frozen teachers supervise these predictions using features extracted from the actual future frames.
The text-grounded target $\mathbf{F}_{\mathrm{tg}}$ comes from the
penultimate-layer patch tokens of SigLIP2~\cite{siglip2}, which provides a language-aligned visual feature space with dense spatial structure.
\rev{These targets are computed from the future images only; the language dependence of $\hat{\mathbf{F}}_{\mathrm{tg}}$ comes from the instruction in the planner input, which is the sense in which we call this branch text-grounded (Sec.~\ref{sec:intro}).}
The spatial-aware target $\mathbf{F}_{\mathrm{spa}}$ comes from the dense
features of Depth Anything~3~\cite{lin2025depth}, which encode object layout,
surface geometry, and relative distances.
We align to these features rather than directly regressing raw depth, so that geometry is represented in a dense-token format before both branches are projected into a common guidance space.
With both targets detached, the planner is optimized to match them:
\begin{equation}
    \revcolor
    \mathcal{L}_{\mathrm{plan}}
    =
    \frac{1}{KV}\sum_{k=1}^{K}\sum_{v=1}^{V}
    \Big[
    \operatorname{MSE}\big(\hat{\mathbf{F}}_{\mathrm{tg}}^{(k,v)},\mathbf{F}_{\mathrm{tg}}^{(k,v)}\big)
    +
    \lambda_{\mathrm{spa}}
    \operatorname{SL_1}\big(\hat{\mathbf{F}}_{\mathrm{spa}}^{(k,v)},\mathbf{F}_{\mathrm{spa}}^{(k,v)}\big)
    \Big],
    \label{eq:planner_loss}
\end{equation}
where $\operatorname{MSE}$ and $\operatorname{SL_1}$ denote the mean-squared and
smooth-$L_1$ losses, both averaged over tokens and feature channels,
$\lambda_{\mathrm{spa}}$ balances the two terms, and \rev{$\mathbf{F}_m^{(k,v)}$ is the teacher feature of view $v$ at the $k$-th future keyframe, so the loss is averaged over the $K$ predicted keyframes and the $V$ camera views}.

\subsection{Semantic-Guided World-Action Model}
\label{sec:method_guidance}

We build the world-action model on a pretrained video diffusion model and couple
its video expert with an action expert through joint
attention~\cite{genieenvisioner,fastwam}.
To condition it on the predicted foresight, we fuse the text-grounded and
spatial-aware maps into a single guidance representation $\mathbf{Z}$ through a
learned gate:
\begin{equation}
    \mathbf{Z}
    =
    \operatorname{Proj}_{\mathrm{tg}}(\hat{\mathbf{F}}_{\mathrm{tg}})
    +
    \sigma(g)\,
    \operatorname{Proj}_{\mathrm{spa}}(\hat{\mathbf{F}}_{\mathrm{spa}}),
    \label{eq:fusion}
\end{equation}
where $\operatorname{Proj}_{\mathrm{tg}}$ and $\operatorname{Proj}_{\mathrm{spa}}$ are
linear projections into the guidance space, $\sigma$ is the sigmoid function, and
$g$ is a learnable scalar gate.
Rather than replacing the original text conditioning, we retain the original text cross-attention and inject $\mathbf{Z}$ into every block of the video expert through an additional parallel cross-attention branch. 
Each guidance token carries a positional encoding of its keyframe time and spatial
location, so it guides the region it describes.
To amplify the effect of the guidance at inference through classifier-free
guidance, we randomly drop the guidance during training so that the model also
learns to generate without it.
The action expert needs no direct injection, since it inherits the guidance
through its joint attention with the video expert.
Both experts are trained with a flow-matching objective.
For each expert, we form $\mathbf{x}_\tau=(1-\tau)\mathbf{y}+\tau\boldsymbol{\epsilon}$
from the clean target $\mathbf{y}$ and Gaussian noise $\boldsymbol{\epsilon}$ at a
sampled time $\tau\in[0,1]$, and regress the velocity
$\mathbf{u}=\boldsymbol{\epsilon}-\mathbf{y}$:
\begin{equation}
    \mathcal{L}_{\mathrm{gen}}
    =
    \underbrace{\mathbb{E}\big\|\hat{\mathbf{v}}_{\mathrm{vid}}-\mathbf{u}_{\mathrm{vid}}\big\|_2^2}_{\mathcal{L}_{\mathrm{vid}}}
    +
    \lambda_{\mathrm{act}}\,
    \underbrace{\mathbb{E}\big\|\hat{\mathbf{v}}_{\mathrm{act}}-\mathbf{u}_{\mathrm{act}}\big\|_2^2}_{\mathcal{L}_{\mathrm{act}}}.
    \label{eq:gen_loss}
\end{equation}
Here $\hat{\mathbf{v}}$ denotes the velocity predicted by each expert, the
subscripts $\mathrm{vid}$ and $\mathrm{act}$ refer to the video and action
experts, and $\lambda_{\mathrm{act}}$ balances the two terms.
\TODO{State the normalization of $\mathbf{y}_{\mathrm{vid}}$ and $\mathbf{y}_{\mathrm{act}}$ that gives $\lambda_{\mathrm{act}}{=}1$ its meaning (e.g., per-channel standardized VAE latents; per-dimension standardized actions).}
% Evidence (line A, HPC3 SG-WAM): pixels -> [-1,1] (ge_act/data/libero_fastwam_hdf5_dataset.py:478); latents (z-mean)/std per channel with LTX VAE
% latents_mean/latents_std, scaling_factor 1.0 (ge_act/utils/data_utils.py:201,242-252). Actions/state: per-suite per-dimension z-score, std+1e-6
% (libero_fastwam_hdf5_dataset.py:237-270,482-485; stats ge_act/configs/ltx_model/libero/libero_fastwam_mix.json {suite}_eef / {suite}_state_eef;
% absolute 7-D EEF actions). Older loader data/libero_dataset.py:326,442 uses q01/q99 min-max to [-1,1]; which one the paper's run used is unconfirmed.

\subsection{Three-Stage Training Paradigm}
\label{sec:method_training}

\paragraph{Stage 1: Semantic Planner Pre-training.}
From only the current frame and the instruction, the planner must reproduce the
dense features that the frozen teachers extract from the $K$ future keyframes.
We optimize the query-token embeddings and the two resampler heads jointly with
the VLM, keeping all teacher encoders frozen,
under the alignment loss $\mathcal{L}_{\mathrm{plan}}$ in
Eq.~\eqref{eq:planner_loss}.
After this stage, the semantic planner can predict text-grounded and
spatial-aware foresight from a single observation.

\paragraph{Stage 2: Semantic-Guided World Model Co-Training.}
We then initialize the world model from a pretrained video diffusion model and
connect it to the semantic planner through the gated fusion and cross-attention
described in Sec.~\ref{sec:method_guidance}, so that training starts from the
intact pretrained behavior.
The semantic planner and the world model are optimized jointly with
\begin{equation}
    \mathcal{L}_{\mathrm{stage2}}
    =
    \mathcal{L}_{\mathrm{vid}}
    +
    \lambda_{\mathrm{plan}}\,\mathcal{L}_{\mathrm{plan}},
    \label{eq:stage2_loss}
\end{equation}
where $\mathcal{L}_{\mathrm{vid}}$ is the video term of Eq.~\eqref{eq:gen_loss} and
$\lambda_{\mathrm{plan}}$ weights the alignment term, which keeps the planner
producing teacher-consistent foresight while adapting to generation.
This stage lets the world model learn to exploit the foresight before gradients from the action loss are introduced.

\paragraph{Stage 3: Semantic-Guided World-Action Co-Training.}
Finally, we attach the action expert and optimize the planner, the world model,
and the action expert together with
\begin{equation}
    \mathcal{L}
    =
    \mathcal{L}_{\mathrm{gen}}
    +
    \lambda_{\mathrm{plan}}\,\mathcal{L}_{\mathrm{plan}}.
    \label{eq:total_loss}
\end{equation}
Crucially, all components are conditioned on the planner's own predictions
rather than the teacher features, which matches deployment.

\subsection{Closed-Loop Deployment}
\label{sec:method_deploy}

At deployment, the teacher encoders are discarded and \ourmethod runs as a
closed-loop visuomotor controller.
At each control cycle, the planner performs a single forward pass that produces
the guidance for all $K$ keyframes and all camera views; conditioned on this
guidance, the video and action experts jointly denoise the future frames and the
action chunk with $10$ Euler steps of the flow-matching sampler, applying the
semantic guidance at every step.
The action chunk is expressed in the joint space of the robot, i.e., the
joint and gripper commands of each arm.
The robot executes only the first $8$ of the $H{=}36$ predicted actions, which
the low-level joint controllers track, after which a new observation is captured
and the policy re-plans in a receding-horizon manner.
The planner and world-action forward passes are thus amortized over eight
control steps, while frequent re-planning lets the policy correct tracking
errors and react to changes in the scene.
